\documentclass[12pt]{article}

\usepackage[paperwidth=200mm,paperheight=112.5mm,margin=11mm,top=10mm]{geometry}
\usepackage{amsmath,amssymb,booktabs,array}
\usepackage[T1]{fontenc}
\usepackage{lmodern}
\usepackage{graphicx}

\pagestyle{empty}
\setlength{\parindent}{0pt}
\setlength{\abovedisplayskip}{9pt}
\setlength{\belowdisplayskip}{9pt}

\newcommand{\R}{\mathbb{R}}
\newcommand{\T}{\mathsf{T}}
\DeclareMathOperator{\col}{col}
\DeclareMathOperator{\range}{range}
\newcommand{\fbase}{f_{\mathrm{base}}}
\newcommand{\ANN}{\mathrm{ANN}}
\newcommand{\thaux}{\theta_{\mathrm{aux}}}
\newcommand{\Phiv}{\Phi_{\bar v}}
\newcommand{\Fann}{F^{\ANN}_{\theta_a}}
\newcommand{\Ftil}{\widetilde F^{\ANN}_{\theta_a}}
\newcommand{\panel}[1]{\subsection*{\Large #1}\vspace{1mm}}

\begin{document}

%======================================================================
\panel{Is the learned component orthogonal to the baseline?}

Lemma~3 asks for $\Phi^{\T}\Ftil=0$. The criterion is therefore
\emph{numerical} zero, not ``small''.

\vspace{1mm}
$\Ftil$ is the component as it is \emph{applied in the loop}, which for our
implementation is not the raw network output: the projection is a correction
subtracted inside the state transition, $x^{+}\!\leftarrow x^{+}-\text{corr}$.

\begin{equation*}
\Ftil=
\begin{cases}
\Fann, & \text{no projection},\\[1mm]
\Fann-\Phiv\thaux, & \text{tangent arm},\\[1mm]
\Fann-\left[\Phiv\thaux+\alpha\,\fbase(\bar v,X_b,U)\right], & \text{affine arm}.
\end{cases}
\end{equation*}

Two overlaps, both scale free. The first is the share of the learned component
that the ten physical parameters $v$ could have produced instead; the second is
its alignment with the baseline's own output, i.e.\ the trained model with the
ANN silenced. For a single column the projection reduces to a cosine.

\begin{equation*}
\rho_{\Phi}=\frac{\bigl\|\Phiv\Phiv^{+}\Ftil\bigr\|}{\bigl\|\Ftil\bigr\|},
\qquad\qquad
\rho_{f}=\frac{\bigl|\langle \fbase,\Ftil\rangle\bigr|}
              {\bigl\|\fbase\bigr\|\,\bigl\|\Ftil\bigr\|}.
\end{equation*}

\newpage
%======================================================================
\panel{Measured overlap, per arm}

\begin{center}
\begin{tabular}{@{}lcc@{}}
\toprule
& $\rho_{\Phi}$ & $\rho_{f}$\\
& \small parameter directions & \small baseline output\\
\midrule
no projection      & $\mathbf{0.87}$      & $1.9\cdot10^{-2}$\\
OBC tangent (10 col) & $8\cdot10^{-8}$    & $1.2\cdot10^{-2}$\\
OBC affine (11 col)  & $2\cdot10^{-7}$    & $1.4\cdot10^{-8}$\\
\bottomrule
\end{tabular}
\end{center}

\vspace{2mm}
Without the projection, $87\%$ of the learned component is reproducible by a
change of the physical parameters. That is the negation mechanism, measured
rather than argued: those are exactly the directions along which $v$ can drift
without the fit noticing.

\vspace{1mm}
The tangent arm removes that share to numerical zero. It leaves
$\rho_{f}=1.2\cdot10^{-2}$, because $\fbase$ does not lie in
$\range(\Phiv)$ once the baseline is nonlinear in its parameters. Removing it
is what the eleventh column of the affine arm is for, and the affine arm does
remove it.

\vspace{1mm}
The projection is not a small correction: it discards about two thirds of the
raw network output, $\|\Fann\|=16.6\to\|\Ftil\|=6.4$ for the tangent arm and
$10.8\to3.8$ for the affine arm.

\vspace{2mm}
\footnotesize
Each arm is evaluated at its own $\bar v$, read from its own checkpoint, over
the $671\,944$ one-step reference tuples the basis is built from. The single
column used for $\rho_{f}$ is $\fbase(\bar v,X_b,U)$, the quantity the question
asks about; slide~4 of 18-09 writes the affine column as
$\Gamma_{\bar v}=\col_k[\fbase-\Phiv\bar v]$. The two differ by a column already
inside $\range(\Phiv)$, so $\widetilde\Phi_{\bar v}$ spans the same subspace and
the projection is identical; only the single-column cosine differs.
Orthogonality is measured on the reference set, which is where it is
constructed; behaviour away from it is not measured here.

\newpage
%======================================================================
\panel{Backup: the same numbers as a figure}

\begin{center}
\includegraphics[height=0.74\textheight]{../figures/orthogonality_applied.pdf}
\end{center}

\end{document}
